% GENERATED FILE. Edit canonical lessons, then run npm run generate:books.
% canonical-source-hash: fnv1a64:f7e255d22429e778
% canonical-lessons: BN-C427-grishmo, BN-C427-shorot, BN-C427-mousum, BN-C427-tapmatra, BN-C427-uttap

\chapter{Seasons and Heat}
\label{ch:bn-seasons-and-heat}
\hlchaptermodality{427}

\begin{chapteropening}
\textbf{By the end of this chapter:} \emph{I can name the hot season and autumn, a fruit's season, and say how hot it is.}

Complete the last lesson of chapter 427: I can name the hot season and autumn, a fruit's season, and say how hot it is.
\end{chapteropening}

\section[\texorpdfstring{\bn{গ্রীষ্ম} (grīshmô)}{গ্রীষ্ম (grīshmô)}]{\bn{গ্রীষ্ম} (grīshmô) — summer; the hot season}

\label{lesson:BN-C427-grishmo}



\begin{quote}
Before the new one: say the Bengali for bitterly cold, then the Bengali for to melt.
\end{quote}

\subsection*{You'll want to know: \bn{গ্রীষ্ম}}
\textbf{\bn{গ্রীষ্ম}} — \emph{grīshmô} — \textquotedblleft{}summer; the hot season\textquotedblright{}.

\begin{quote}
\bn{গ্রীষ্মে} \bn{দুপুরে} \bn{রাস্তা} \bn{ফাঁকা} \bn{থাকে।}
\end{quote}

\begin{quote}
\emph{grīshme dupure rāstā phā̃kā thāke.}
\end{quote}

\begin{quote}
In summer the streets are empty at midday.
\end{quote}

The calendar word; in talk people say \bn{গরম}, the heat.

\begin{cousinweb}
\textbf{\bn{গ্রীষ্ম}} (\emph{grīshmô}) is Sanskrit \emph{grīṣma}, \textquotedblleft{}summer, the hot season\textquotedblright{} (Monier-Williams), taken over whole.
\end{cousinweb}

\begin{grammarlens}[title={What you've built}]
One more word for getting around.
\end{grammarlens}

\subsection*{Your turn}
\begin{itemize}
\raggedright
  \item \emph{Say it:} \emph{grīshmô}
  \item \emph{Say it:} \emph{grīshmô}, once more
  \item \emph{Say it:} \emph{grīshme dupure rāstā phā̃kā thāke.} — In summer the streets are empty at midday.
  \item \emph{Recall:} say the Bengali for bitterly cold, then the Bengali for to melt, then say \emph{grīshmô} again
\end{itemize}

\begin{tcolorbox}[breakable,colback=teal!4,colframe=teal!35!black,title={Before you move on}]
What does \bn{গ্রীষ্ম} mean? (\textquotedblleft{}Summer; the hot season\textquotedblright{}.) Say it once more.
\end{tcolorbox}
\section[\texorpdfstring{\bn{শরৎ} (shôrôt)}{শরৎ (shôrôt)}]{\bn{শরৎ} (shôrôt) — autumn (the season after the rains)}

\label{lesson:BN-C427-shorot}



\begin{quote}
Before the new one: say the Bengali for to melt, then the Bengali for summer.
\end{quote}

\subsection*{You'll want to know: \bn{শরৎ}}
\textbf{\bn{শরৎ}} — \emph{shôrôt} — \textquotedblleft{}autumn (the season after the rains)\textquotedblright{}.

\begin{quote}
\bn{শরতে} \bn{আকাশ} \bn{নীল} \bn{আর} \bn{পরিষ্কার} \bn{থাকে।}
\end{quote}

\begin{quote}
\emph{shôrôte ākāsh nīl ār pôrishkār thāke.}
\end{quote}

\begin{quote}
In autumn the sky stays blue and clear.
\end{quote}

Before an ending \bn{ৎ} is written \bn{ত}: \bn{শরতে}, in autumn. It is the season of the Durga Puja.

\begin{cousinweb}
\textbf{\bn{শরৎ}} (\emph{shôrôt}) is Sanskrit \emph{śarad}, \textquotedblleft{}autumn\textquotedblright{} (Monier-Williams), the season of ripening after the rains.
\end{cousinweb}

\begin{grammarlens}[title={What you've built}]
One more word for getting around.
\end{grammarlens}

\subsection*{Your turn}
\begin{itemize}
\raggedright
  \item \emph{Say it:} \emph{shôrôt}
  \item \emph{Say it:} \emph{shôrôt}, once more
  \item \emph{Say it:} \emph{shôrôte ākāsh nīl ār pôrishkār thāke.} — In autumn the sky stays blue and clear.
  \item \emph{Recall:} say the Bengali for to melt, then the Bengali for summer, then say \emph{shôrôt} again
\end{itemize}

\begin{tcolorbox}[breakable,colback=teal!4,colframe=teal!35!black,title={Before you move on}]
What does \bn{শরৎ} mean? (\textquotedblleft{}Autumn (the season after the rains)\textquotedblright{}.) Say it once more.
\end{tcolorbox}
\section[\texorpdfstring{\bn{মৌসুম} (mousum)}{মৌসুম (mousum)}]{\bn{মৌসুম} (mousum) — a season (for a crop or an activity)}

\label{lesson:BN-C427-mousum}



\begin{quote}
Before the new one: say the Bengali for summer, then the Bengali for autumn.
\end{quote}

\subsection*{You'll want to know: \bn{মৌসুম}}
\textbf{\bn{মৌসুম}} — \emph{mousum} — \textquotedblleft{}a season (for a crop or an activity)\textquotedblright{}.

\begin{quote}
\bn{আমের} \bn{মৌসুমে} \bn{বাজার} \bn{আমে} \bn{ভরা।}
\end{quote}

\begin{quote}
\emph{āmer mousume bājār āme bhôrā.}
\end{quote}

\begin{quote}
In mango season the market is full of mangoes.
\end{quote}

The English word monsoon comes from the same Arabic source.

\begin{cousinweb}
\textbf{\bn{মৌসুম}} (\emph{mousum}) is usually traced to Arabic \emph{mawsim}, \textquotedblleft{}a season, a time of year\textquotedblright{}, through Persian; Portuguese sailors took the word for the trade winds, and it reached English as \emph{monsoon}.
\end{cousinweb}

\begin{grammarlens}[title={What you've built}]
One more word for getting around.
\end{grammarlens}

\subsection*{Your turn}
\begin{itemize}
\raggedright
  \item \emph{Say it:} \emph{mousum}
  \item \emph{Say it:} \emph{mousum}, once more
  \item \emph{Say it:} \emph{āmer mousume bājār āme bhôrā.} — In mango season the market is full of mangoes.
  \item \emph{Recall:} say the Bengali for summer, then the Bengali for autumn, then say \emph{mousum} again
\end{itemize}

\begin{tcolorbox}[breakable,colback=teal!4,colframe=teal!35!black,title={Before you move on}]
What does \bn{মৌসুম} mean? (\textquotedblleft{}A season (for a crop or an activity)\textquotedblright{}.) Say it once more.
\end{tcolorbox}
\section[\texorpdfstring{\bn{তাপমাত্রা} (tāpmātrā)}{তাপমাত্রা (tāpmātrā)}]{\bn{তাপমাত্রা} (tāpmātrā) — temperature}

\label{lesson:BN-C427-tapmatra}



\begin{quote}
Before the new one: say the Bengali for autumn, then the Bengali for a season (for a crop or an activity).
\end{quote}

\subsection*{You'll want to know: \bn{তাপমাত্রা}}
\textbf{\bn{তাপমাত্রা}} — \emph{tāpmātrā} — \textquotedblleft{}temperature\textquotedblright{}.

\begin{quote}
\bn{আজ} \bn{শহরের} \bn{তাপমাত্রা} \bn{চল্লিশ} \bn{ডিগ্রি।}
\end{quote}

\begin{quote}
\emph{āj shôhôrer tāpmātrā chôllish ḍigri.}
\end{quote}

\begin{quote}
Today the temperature in the city is forty degrees.
\end{quote}

\begin{cousinweb}
\textbf{\bn{তাপমাত্রা}} (\emph{tāpmātrā}) joins Sanskrit \emph{tāpa}, \textquotedblleft{}heat\textquotedblright{}, and \emph{mātrā}, \textquotedblleft{}a measure\textquotedblright{} (Monier-Williams for both): the measure of heat, a modern coinage.
\end{cousinweb}

\begin{grammarlens}[title={What you've built}]
One more word for getting around.
\end{grammarlens}

\subsection*{Your turn}
\begin{itemize}
\raggedright
  \item \emph{Say it:} \emph{tāpmātrā}
  \item \emph{Say it:} \emph{tāpmātrā}, once more
  \item \emph{Say it:} \emph{āj shôhôrer tāpmātrā chôllish ḍigri.} — Today the temperature in the city is forty degrees.
  \item \emph{Recall:} say the Bengali for autumn, then the Bengali for a season (for a crop or an activity), then say \emph{tāpmātrā} again
\end{itemize}

\begin{tcolorbox}[breakable,colback=teal!4,colframe=teal!35!black,title={Before you move on}]
What does \bn{তাপমাত্রা} mean? (\textquotedblleft{}Temperature\textquotedblright{}.) Say it once more.
\end{tcolorbox}
\section[\texorpdfstring{\bn{উত্তাপ} (uttāp)}{উত্তাপ (uttāp)}]{\bn{উত্তাপ} (uttāp) — heat; warmth}

\label{lesson:BN-C427-uttap}



\begin{quote}
Before the new one: say the Bengali for a season (for a crop or an activity), then the Bengali for temperature.
\end{quote}

\subsection*{You'll want to know: \bn{উত্তাপ}}
\textbf{\bn{উত্তাপ}} — \emph{uttāp} — \textquotedblleft{}heat; warmth\textquotedblright{}.

\begin{quote}
\bn{আগুনের} \bn{উত্তাপে} \bn{ঘর} \bn{গরম} \bn{হল।}
\end{quote}

\begin{quote}
\emph{āguner uttāpe ghôr gôrôm hôlô.}
\end{quote}

\begin{quote}
The heat of the fire warmed the room.
\end{quote}

\begin{cousinweb}
\textbf{\bn{উত্তাপ}} (\emph{uttāp}) is Sanskrit \emph{ut-tāpa}, \textquotedblleft{}great heat, glow\textquotedblright{} (Monier-Williams), from \emph{ut}, up, and \emph{tāpa}, heat.
\end{cousinweb}

\begin{grammarlens}[title={What you've built}]
One more word for getting around.
\end{grammarlens}

\subsection*{Your turn}
\begin{itemize}
\raggedright
  \item \emph{Say it:} \emph{uttāp}
  \item \emph{Say it:} \emph{uttāp}, once more
  \item \emph{Say it:} \emph{āguner uttāpe ghôr gôrôm hôlô.} — The heat of the fire warmed the room.
  \item \emph{Recall:} say the Bengali for a season (for a crop or an activity), then the Bengali for temperature, then say \emph{uttāp} again
\end{itemize}

\begin{tcolorbox}[breakable,colback=teal!4,colframe=teal!35!black,title={Before you move on}]
What does \bn{উত্তাপ} mean? (\textquotedblleft{}Heat; warmth\textquotedblright{}.) Say it once more.
\end{tcolorbox}
